\begin{circuitikz}[
    american,
    voltage dir=RP,
    >=latex,
    %>={Latex[length=5mm, width=2mm]},
    alt={Diagram showing a single-rung ladder logic diagram featuring a Retentive Timer On (RTO) instruction labeled AutoTimer. On the left rail, a normally open contact Auto_Mode connects to the input of the RTO instruction block. The timer is configured with a Preset (PRE) value of 28800 and an Accumulated (ACC) value of 28804. Green power flow highlighting extends from the left vertical rail through the Auto_Mode contact and into both the enable (EN) and done (DN) states of the timer block, indicating that the timer is enabled, has accumulated past its preset value, and has set its done bit to true.}
    ]

    % Define the number of rungs in the diagram
    \def\numRungs{1}

    % Calculate the length of the vertical power rails dynamically based on number of rungs
    \pgfmathsetmacro{\railLength}{-2 * \numRungs}

    %======================================================================================================
    % Component Grid
    % Spacing = 3x, 2y
    %======================================================================================================
    \foreach \i in {0,...,5} {
        \foreach \j in {0,...,20} { % Change to 0,...,19 if you need 20 individual Y points down to -40
        
        % Calculate spatial positions
        \pgfmathsetmacro{\px}{\i * 3}
        \pgfmathsetmacro{\py}{-1 - (\j * 2)}
        
        % Name coordinate using loop indices: (C-column-row) e.g., (C-0-0), (C-5-10)
        \coordinate (C-\i-\j) at (\px, \py);
        
        % OPTIONAL: Uncomment the 2 lines below to display grid labels while building
        % \fill[red] (C-\i-\j) circle (1.5pt);
        %\node[anchor=south west, scale=0.5, gray] at (C-\i-\j) {(\i,\j)};
        }
    }

    %======================================================================================================
    % Foreground Layer
    % Only contains the vertical power rails
    % Length of the rails is automatic, calculated based on the number of rungs defined above
    %======================================================================================================
    \begin{pgfonlayer}{ft}

        %======================================================================================================
        % Scope sets options for both lines - BakerBlack, very thick
        %======================================================================================================
        \begin{scope}[
            draw=BakerBlack,
            text=BakerBlack,
            color=BakerBlack,
            font=\small,
            ultra thick
        ]
        
            \draw (0,0) -- (0,\railLength);

            \draw (15,0) -- (15,\railLength);

        \end{scope}
    \end{pgfonlayer}

    %======================================================================================================
    % Main Layer
    % All actual logic goes here
    %======================================================================================================
    \begin{pgfonlayer}{main}
        
        %======================================================================================================
        % Scope sets options for both lines - BakerBlack for draw and text
        %======================================================================================================
        \begin{scope}[
            draw=BakerBlack,
            text=BakerBlack,
            color=BakerBlack,
            font=\small
        ]

            %=================================================================
            % Rung 1
            %=================================================================
            \draw (C-0-0) to[C, l={Auto\_Mode}] (C-1-0) -- (C-3-0) pic (T1) {TON};
            \draw (T1-OUT) -- (C-5-0);

            % Instruction box labels
            \node[anchor=south] at (T1-TOP) {RTO};        % Instruction title
            \node[anchor=north] at (T1-TOP) {AutoTimer};  % Instruction variable
            \node[anchor=west] at (T1-PRE) {28800};       % PRE value
            \node[anchor=west] at (T1-ACC) {28804};       % ACC value

        \end{scope}
    \end{pgfonlayer}

    %=================================================================
    % Background layer
    %=================================================================
    \begin{pgfonlayer}{bg}
        
        %==================================================================
        % Highlight Scope
        % This layer exists for highlighting true logic, the scope sets the
        % highlight options without having to repeat them every \draw
        % draw=UpNorthGreen, line width=10pt, opacity=0.4
        %==================================================================
        \begin{scope}[
            draw=UpNorthGreen,
            line width=10pt,
            opacity=0.4
        ]

            %==================================================================
            % Left vertical rung - leave this alone, it should ALWAYS be green
            % It automatically draws itself to the length defined at the top
            %==================================================================
            \draw (0,0) -- (0,\railLength);
            
            %=================================================================
            % Rung 1
            %=================================================================
            \draw (C-0-0) -- (C-1-0);

            \draw (T1-EN1) -- (T1-EN2);
            \draw (T1-DN1) -- (T1-DN2);

        \end{scope}
    \end{pgfonlayer}
    
\end{circuitikz}